%
% Chapitre "Structure d'un rapport technique"
%

\chapter{Concept retenu}
\label{s:concept_retenu}
\section{Matrice de décision}
Texte

\section{Analyse de la matrice décisionnelle}
Texte

\begin{table}[h!]
\centering
\caption{Matrice décisionnelle}
\label{tab:matricede}
\small
\begin{tabular}{|l|c||c|c|c|}
\hline
Critères d'évaluation & Pond. & Concept 1 & Concept 2 & Concept 3 \\
\hline\hline
\textbf{\ref{Autonomie} Fonctionnement autonome du système} & \textbf{30\%} & & &  \\
\hline
\ref{robutesse_autonomie} Marge de l'état d'autonomie & 10\% & 2,9\% & 7,9\% & 10\% \\
\ref{consommation_electrique} Consommation énergétique & 8\% & 5\% & 0\% & 4,25\% \\
\ref{resilience_systeme} Proportion des fonctions protégées & 6\% & 3,6\% & 2,4\% & 4,8\% \\
\hline

\textbf{\ref{Robustesse} Résistance et entretien du système} & \textbf{30\%} & & &  \\
\hline
\ref{conditions_climatiques} Plage opérationnelle de température & 10\% & 0\% & 0\% & -30 \\
\ref{installation_deploiement} Encombrement et complexité d'installation & 6\% & 4,8\% & ?? & ?? \\
\ref{resistance_neige} Résistance mécanique à la neige & 12\% & 12\% & 12\% & 12\% \\
\hline

\textbf{\ref{donnees} Observation et gestion de données} & \textbf{25\%} & & & \\
\hline
\ref{detection_mammiferes} Couverture de la zone d'observation & 9\% & 9\% & 9\% & 93 \\
\ref{conservation_donnees} Capacité de stockage & 5\% & 0\% & 5\% & 134 \\
\ref{performance_communication} Performance de communication distante & 6\% & 5,1\% & ?? & ?? \\
\ref{exploitation_scientifique} Qualité scientifique des données & 5\% & 3,75\% & ?? & ?? \\
\hline

\textbf{\ref{impacts_couts} Réduire impacts et coûts} & \textbf{15\%} & & &  \\
\hline
\ref{impact_lumineux} Longueur d'onde de l'éclairage & 3\% & 3\% & 2,07\% & ?? \\
\ref{impact_environnemental} Nombre de remplacements de batterie par an & 2\% & 1,5\% & 2 & 1 \\
\ref{couts_transport} Coûts d'installation et de transport & 4\% & 3,1\% & 2000 & 2200 \\
\ref{couts_entretien} Coûts d'opération et d'entretien & 3\% & 1,97\% & 1500 & 1800 \\
\ref{couts_fabrication} Coûts de fabrication & 3\% & 2,58\% & 1500 & 2169 \\
\hline\hline

\textbf{Total} & \textbf{100\%} & & & \\
\hline
\end{tabular}
\end{table}

\section{Description du concept retenu}
\textbf{Sopporter les conditions métérologiques}
\begin{itemize}
    \item Protection du système avec un boitier en bois
\end{itemize}

\textbf{Détection d'un mouvement et enregistrement de vidéos}
\begin{itemize}
    \item La détection de mouvement sera assurée par des plaques piézoélectriques
    \item Les vidéos seront capturées par l'ArduCam IMX219
    \item Les données seront enregistrées sur une carte MicroSD(HC)
\end{itemize}

\textbf{Interface utilisateur}
\begin{itemize}
    \item Le système pourra être accédé par l'intermédiare du Tunnel Cloudflare
    \item La communication entre le système et l'extérieur sera assurée par l'Iridium 9603 SBD
\end{itemize}

\textbf{Gestion de l'énergie}
\begin{itemize}
    \item Le système sera alimenté par un blocs de batteries au Lithium-thionyl chloride
\end{itemize}

\subsection{Gestion de l'énergie}
L'entièreté du système sera alimenté par le bloc de batteries au Lithium-thionyl chloride.
Ces batteries hautes densistés combinées avec un système qui consomme au maximum 11.3 Watts pour
de très courtes pédiodes lors d'enregistrement vidéos, assurera une autonomie d'au moins 12 mois
pour le système.

Si l'on programme un système qui s'active uniquement lorsqu'un signal est envoyé par la plaque 
de pression, on peut estimer une consommation maximale en veille de 0,1 W, l'on garde la 
plaque piézoélectrique active. Avec une consommation 

\section{Conclusion}
Texte